\section*{Implementation Details}
\emph{Instances and state graph.} With the depot horizon $H=l_0$, the six static node features are $x_i/100$, $y_i/100$,
$d_i/Q$, $e_i/H$, $l_i/H$ and $\sigma_i/H$. The six dynamic features are computed on the \emph{current} solution $x$
(not on the best one): waiting time, slack $l_i-b_i$ and forward time slack (Savelsbergh, 1992), each divided by $H$;
cumulative load divided by $Q$; relative position in the route; and a routed flag (the depot has zeros). The edge set
contains, for every customer, directed arcs to its 10 nearest customers, arcs between the depot and its 10 nearest
customers, and both directions of every arc of the current routes (depot legs included). Each edge carries
$c_{ij}/c_{\max}$ and a flag (0: neighbour arc, 1: route arc). The graph therefore grows with $n$, and mean and max
pooling give an embedding of fixed size. The global vector holds $t/T$, $\tau/\tau_0$, the number of iterations since
the last new best divided by $0.2T$ (capped at 1), $(f(x)-f(x^\star))/f(x^\star)$, $(f(x_0)-f(x^\star))/f(x_0)$ with
$x_0$ the solution after vehicle reduction, the number of routes of $x$ divided by that of the I1 solution, and the
success rate (share of uses that improved the current solution) of each of the nine
operators in the last segment. \sgc{} appends $n/200$ and $T/5{,}000$.

\emph{Roulette wheel.} After each segment, $w_o\leftarrow\max\{(1-\rho)w_o+\rho\,\psi_o/u_o,\ 10^{-3}\}$, with
$\psi_o/u_o:=0$ for an operator that was not used, so the weight of an unused operator decays geometrically towards the
floor. Destroy and repair probabilities are the weights normalised within each group. Scores per use are 33 (new
best), 9 (better than the current solution) and 13 (accepted worse solution). A repair that cannot insert every customer is
rejected and scores 0.

\emph{Factorised policy.} The three heads are categorical distributions over $a^1$ (19), $a^2$ (5) and $a^3$ (3)
computed from the same embedding, i.e.\ conditionally independent given the state:
$\log\pi(a\mid s)=\sum_k\log\pi_k(a^k\mid s)$, and the entropy used in the PPO loss is the sum of the three head
entropies (exact for this factorisation). In an ablation, the removed head is masked to its expert value and
contributes zero to both. The output weights are scaled by 0.1 at initialisation and the biases of the expert values
are set to $3$ ($a^1$, RW) and $2$ ($a^2$, default range; $a^3$, $m=1$); the neutral-prior variant sets all biases to 0.

\emph{PPO.} Rewards are divided by a running standard deviation; each update runs 4 epochs over minibatches of 256
decisions, with value coefficient 0.5 and gradient-norm clipping at 0.5. Every 5 updates the greedy policy is scored
against \Rw{} on 24 fixed generated instances ($n=100$); this validation is only logged, and the \emph{final}
checkpoint is evaluated. Training instances, validation instances (seed 12345) and the 24 tuning instances of the
tuned baselines (seeds 90000--90023) are drawn from disjoint random streams of Solomon's generator; no benchmark
instance is used for any selection.

\emph{Objective and reward.} The best solution $x^\star$ is always tracked with $f=10^4\mathrm{NV}+\mathrm{TD}$. Over
all \qLexRuns{} runs of this study, the largest distance difference between two solutions of the same instance is
\qLexRange{} (the largest TD is \qLexTdMax{}), so $f$ orders every pair of our solutions lexicographically (NV first),
also at $n=400$. The reward of Eq.~(1) is a learning signal only: one saved vehicle is worth $\lambda=5$ points,
i.e.\ as much as a 5\% reduction of the initial distance. Unlike $f$, the reward therefore does not enforce the
vehicle-first order: it is a shaping choice, and Table~\ref{tab:reward} varies $\lambda$.
